\section{Persistence and one fixed smooth metric}\label{sec:limit}

\subsection{A finite norm suffices to preserve earlier signs}

\begin{lemma}\label{lem:persistence}
Let a smooth positive pair \(\mathcal P\) have a simple positive
eigenvalue \(\lambda\) and a real eigenfunction with a sign
certificate \(\mathcal C\). For any open interval
\(I\subset(0,\infty)\) containing \(\lambda\), there are
an integer \(s\ge1\) and \(\varepsilon>0\) such that every
smooth positive pair \(\mathcal P'\) with
\(\|\mathcal P'-\mathcal P\|_{C^s}\le\varepsilon\)
has a simple eigenvalue in \(I\) and a real eigenfunction
with the same certificate disks.
\end{lemma}

\begin{proof}
We specify the compactness and regularity facts used both here
and in Lemma~\ref{lem:simple}.
On \(H^1(B)\), the shifted form
\[
 a(f,v)=\int_B A(\dd f,\dd v)\dd\mu+
                     \int_B\rho fv\dd\mu
\]
is continuous, symmetric and coercive.
The inclusion \(H^1(B)\hookrightarrow L^2(\rho\dd\mu)\)
is continuous, compact and dense.
These properties follow in a finite atlas from the local Sobolev
compactness theorem \cite[Theorem~3.49, p.~76]{Hunter};
positive smooth densities give equivalent \(L^2\) norms.
Laugesen's discrete spectral theorem
\cite[Theorem~4.1, pp.~28--32]{Laugesen}, with form space
\(H^1\) and Hilbert space \(L^2(\rho\dd\mu)\), gives an orthonormal
eigenbasis for the shifted problem with eigenvalues tending to
infinity. Subtracting one gives the nonnegative discrete spectrum
of \(-L\), with finite multiplicities.
All eigenfunctions can be chosen real.

For completeness, weak eigenfunctions are smooth.
In a coordinate chart write \(\dd\mu=m(x)\dd x\).
The divergence coefficients are \(mA^{ab}\), and the right-hand
side is a smooth density times the eigenfunction.
Local uniform ellipticity and
\cite[Theorems~4.27--4.28, pp.~112--114]{Hunter}
first give \(H^2_{\rm loc}\) and then successively every finite
Sobolev order. Their hypotheses on coefficient derivatives and
the Sobolev regularity of the right-hand side are met at each
step. Sobolev embedding and a finite cover give global smoothness.
The same argument applies to a weak limit solving a fixed smooth
eigenfunction equation.

The eigenbasis also gives the min--max characterization
\begin{equation}\label{eq:minmax}
 \lambda_n=\inf_{\dim F=n+1}\ \sup_{0\ne f\in F}
 \frac{\int_B A(\dd f,\dd f)\dd\mu}
      {\int_B\rho f^2\dd\mu},\qquad n=0,1,\ldots .
\end{equation}
Indeed, the first \(n+1\) eigenvectors give the upper bound.
Every \((n+1)\)-dimensional \(F\) contains a nonzero vector
orthogonal to the first \(n\), whose eigenbasis expansion gives
the lower bound. The characterization is on the common form
domain \(H^1\), including when the density varies.
Uniform multiplicative comparisons of nearby positive tensors
and densities compare all the quotients in \eqref{eq:minmax}.
Thus \(C^0\) convergence of pairs implies convergence of every
fixed ordered eigenvalue. A smaller interval with closure in \(I\)
can separate the simple \(\lambda\) from its neighboring
ordered eigenvalues. For all sufficiently close pairs it contains
exactly the corresponding eigenvalue, which is still simple.

We need convergence of its normalized eigenfunction in \(C^0\).
The same interior estimates, localization by a finite smooth
partition of unity and absorption of lower-order terms give,
for the fixed smooth operator \(L\),
\begin{equation}\label{eq:elliptic}
 \|f\|_{H^{a+2}}\le C_a
       \bigl(\|Lf\|_{H^a}+\|f\|_{L^2}\bigr),\qquad a\ge0.
\end{equation}
One way to see the absorption is to apply the interior estimates
on a slightly larger chart to cutoff multiples of \(f\).
The commutator with the cutoff is first order.
The inequalities
\(\|f\|_{H^{a+1}}\le
 \epsilon\|f\|_{H^{a+2}}+C_\epsilon\|f\|_{L^2}\)
follow from the elementary Fourier weight inequality in each
chart. A finite sum and sufficiently small \(\epsilon\) absorb
those commutators and give \eqref{eq:elliptic}.

A sufficiently small \(C^4\) change of the pair makes
\(L'-L\) small as a map \(H^{a+2}\to H^a\) for \(a=0,2\).
In coordinates its principal coefficients have no coefficient
derivatives, and its first-order coefficients have at most one;
the stated finite norm bounds every product-rule coefficient
needed at these two orders. Absorb
\(\|(L'-L)f\|_{H^a}\) into \eqref{eq:elliptic}.
For eigenvalues in the chosen bounded interval, the estimates
at \(a=0\) and then \(a=2\) give uniform \(H^4\) bounds for
weighted-\(L^2\)-normalized eigenfunctions.

If \(\mathcal P_m\to\mathcal P\) in this finite norm, choose
real eigenfunctions \(f_m\) with
\(\int_B\rho_m f_m^2\dd\mu=1\).
A subsequence converges strongly in \(H^3\), and hence in \(C^0\);
the dimension-four embedding \(H^3\hookrightarrow C^0\)
again follows from Hunter's Theorem~3.49 in a finite atlas.
The weak equation passes to the limit, whose norm is one
because \(\rho_m\to\rho\) uniformly. Simplicity identifies the
limit with the given normalized eigenfunction, up to sign.
On the finite collection of compact disks its strict signs
have a positive absolute margin. Uniform convergence retains
all those signs, possibly after a global sign change.
If no sufficiently small finite-norm neighborhood retained
them, a sequence contradicting this conclusion could be
chosen and the same compactness argument would apply.
Finally shrink the radius to put a closed finite-norm ball
inside that neighborhood.
\end{proof}

\subsection{Summable changes and separated exact eigenvalues}

\begin{proof}[Proof of Theorem~\ref{thm:main}]
Take monotone coordinate norms as in Section~\ref{sec:base}.
Choose \(b>0\) such that the closed \(C^1\) ball of radius \(b\)
about \(\mathcal P_0=(g_0^{-1},1)\) consists of positive pairs
inside \(\mathcal U\). Set
\[
 l^{(j)}=\frac14+2^{-j-2}\quad(j\ge0),\qquad
 \epsilon_0=\gamma/4>0 .
\]
We construct \(\mathcal P_j=(A_j,\rho_j)\), round for
\(t<l^{(j)}\), with a simple eigenvalue \(k_j^2\) and
certificate \(\mathcal C_j\).

After stage \(i\), apply Lemma~\ref{lem:persistence} with interval
\((k_i^2/2,2k_i^2)\), and call its norm order and radius
\(s_i,\varepsilon_i\). At stage \(j\), only earlier such
numbers enter the requirements:
\begin{equation}\label{eq:budgets}
 \begin{split}
 r_j&=\max\{j+1,1,s_1,\ldots,s_{j-1}\},\\
 \eta_j&=\min\left\{2^{-j},\,b2^{-j-1},\,
              \min_{1\le i<j}
                    2^{-(j-i+1)}\varepsilon_i\right\}.
 \end{split}
\end{equation}
For \(j=1\) omit the inner minimum. Apply
Proposition~\ref{prop:step} to the fixed input
\(\mathcal P_{j-1}\), collar
\((l^{(j)},l^{(j-1)})\), order \(r_j\) and tolerance \(\eta_j\).
Its positive constant \(c_j\) is known before choosing \(k_j\).
Choose an allowed frequency beyond its threshold satisfying
\[
 k_j>2k_{j-1}\quad(j>1),\qquad c_jk_j^{\gamma/2}>j.
\]
The resulting certificate then obeys
\begin{equation}\label{eq:stagesize}
 |\mathcal C_j|>j k_j^{1+\gamma/2}.
\end{equation}
Use the disks constructed in Lemma~\ref{lem:signs}; their centers
are in \(V\Subset\Omega\), and their radii are at most \(k_j^{-1}\).
Choose also the first frequency large enough to keep all cylinders
in one fixed compact subregion of \(\Omega\).
Only now select \(s_j,\varepsilon_j\) for this output.
They restrict later stages and do not constrain the frequency
that produced them.

For every fixed \(r\), eventually \(r_j\ge r\), and
\(\|\mathcal P_j-\mathcal P_{j-1}\|_{C^r}<2^{-j}\).
The sequence therefore converges in every \(C^r\) to a single
smooth pair \(\mathcal P_\infty=(A_\infty,\rho_\infty)\).
Its total \(C^1\) movement from the round pair is at most
\(b\sum_{j\ge1}2^{-j-1}=b/2\), so it is positive and lies
inside \(\mathcal U\).
For each \(i\), \eqref{eq:budgets} gives
\[
 \|\mathcal P_\infty-\mathcal P_i\|_{C^{s_i}}
 \le\sum_{j>i}2^{-(j-i+1)}\varepsilon_i
       =\varepsilon_i/2 .
\]
Thus Lemma~\ref{lem:persistence} applies simultaneously to every
earlier certificate. There are nonzero real smooth \(v_i\) for
the same limiting operator satisfying
\begin{equation}\label{eq:limitlambda}
 -L_\infty v_i=\lambda_i v_i,\qquad
 k_i^2/2<\lambda_i<2k_i^2,
\end{equation}
with certificate disks \(\mathcal C_i\).
The eigenvalue intervals are disjoint and tend to infinity
because \(k_{i+1}>2k_i\).

\subsection*{An ordinary metric with the prescribed base operator}

Put
\[
 G_\infty^{-1}=A_\infty/\rho_\infty,\qquad
 q=\frac{\rho_\infty\,\dd\mu}{\dd\operatorname{vol}_{G_\infty}},
 \qquad
 g=G_\infty+q^2\dd\varphi^2\quad\hbox{on }B\times S^1 .
\]
The ratio \(q\) of smooth positive densities is smooth and positive.
Hence \(g\) is a smooth positive metric on a closed connected
five-dimensional manifold, fixed independently of \(i\).
Its volume density is
\[
 \dd\operatorname{vol}_g
      =q\,\dd\operatorname{vol}_{G_\infty}\dd\varphi
      =\rho_\infty\,\dd\mu\dd\varphi .
\]
For any function \(f\) independent of \(\varphi\), the coordinate
divergence formula, with \(\mu=m(x)\dd x\), gives
\begin{equation}\label{eq:circle}
 \Delta_g f=\frac{1}{\rho_\infty m}
  \partial_a(\rho_\infty mG_\infty^{ab}\partial_b f)
     =\rho_\infty^{-1}\Div_\mu(A_\infty\dd f)=L_\infty f .
\end{equation}
The pullbacks \(u_i\) of \(v_i\) are therefore exact real
Laplace eigenfunctions on this single metric, at \(\lambda_i\).

\subsection*{An intrinsic Hausdorff measure lower bound}

Fix a closed interval \(I_0\) of positive length in a circle
coordinate chart. For one certificate cylinder \(C_i^{(a)}\),
let \(D_i^{(a)}\) denote its transverse disk.
Continuity and opposite endpoint signs force a zero on each
joining segment. Consequently
\[
 \Z(u_i)\cap(C_i^{(a)}\times I_0)
\]
projects onto \(D_i^{(a)}\times I_0\) by the orthogonal
base-coordinate projection and the circle coordinate.

All these cylinders lie in a single compact coordinate region.
Extend its coordinate functions, and the circle coordinate on
\(I_0\), smoothly with cutoffs to the compact manifold.
Their gradients have fixed bounds in \(g\). Integrating along
curves bounds their differences by a fixed constant times
the intrinsic distance. Composing with any orthogonal projection
gives one Lipschitz constant \(K\), uniform in every disk's
radius and orientation.
The Hausdorff covering inequality therefore implies
\begin{equation}\label{eq:project}
 \HH_g^4\bigl(\Z(u_i)\cap(C_i^{(a)}\times I_0)\bigr)
 \ge K^{-4}\HH_{\mathrm{Euc}}^4(D_i^{(a)}\times I_0)
 \ge c\,\HH_{\mathrm{Euc}}^3(D_i^{(a)})|I_0|.
\end{equation}
The last constant is purely dimensional: the product disk has
four-dimensional Lebesgue volume equal to the disk's
three-volume times \(|I_0|\), while a set of diameter \(d\)
in its four-dimensional affine span has volume at most \(Cd^4\).
Apply this estimate to every countable cover and take the
Hausdorff infimum. The identical three-dimensional covering
normalization identifies the disk's three-volume with a fixed
multiple of \(\HH^3_{\mathrm{Euc}}\).
Thus no convention on the normalization of Hausdorff measure
affects the bound.

The finitely many closed cylinders are in disjoint open balls,
so their compact nodal subsets are pairwise positively separated.
Hausdorff measure is additive on these separated sets directly
from covers of sufficiently small diameter. Summing
\eqref{eq:project} and using \eqref{eq:stagesize} yields
\begin{equation}\label{eq:finalmeasure}
 \HH_g^4(\Z(u_i))\ge c|\mathcal C_i|
                   >c\,i k_i^{1+\gamma/2}.
\end{equation}
Only continuity and strict signs were used; no regularity of
the nodal set is needed.
Finally, \eqref{eq:limitlambda} and \(\epsilon_0=\gamma/4\) give
\[
 \lambda_i^{1/2+\epsilon_0}
 \le 2^{1/2+\epsilon_0}k_i^{1+\gamma/2},\qquad
 \frac{\HH_g^4(\Z(u_i))}{\lambda_i^{1/2+\epsilon_0}}
       \ge c\,2^{-1/2-\epsilon_0}i\longrightarrow\infty .
\]
The metric and positive exponent are fixed for the full sequence.
This proves Theorem~\ref{thm:main}.
\end{proof}
